\scanpage{6}
\mspage{5}
the form
\[
 \deleted{\unsure{X}}A_J=\bigcup_{i\in J}A_i^\circ={\unsure{\varphi|\phi}}^{-1}(J)
 \tag*{(1)}
\]
where $\deleted{\unsure{A}}J$ is an open subset of $I$ (or of $\unsure{\Phi|\Psi|\mathscr{L}}\simeq I$) (\deleted{\unsure{of}} course, this time in $\displaystyle\bigcup_{i\in J}$ there is no question of replacing $A_i^\circ$ by $A_i$ !). The locally closed subsets are of \ensuremath{=} the
\[
 \deleted{\unsure{X}}A_J=\varphi^{-1}(J)=\bigcup_{i\in J}A_i^\circ
\]
where this time $J$ is a locally closed subset of $I$, i.e. such that
\[
 i\leq j\leq k,\quad i,k\in J\quad\Rightarrow\quad j\in J
 \tag*{$(*)$}
\]
(the $\deleted{\unsure{J|I}}I_{\geq i}$, $i\in J$, are open subsets of $I$ covering $J$, and \deleted{\unsure{in this}}\inserted{for such an} open subset, \deleted{$J$}\inserted{$J\cap I_{\geq i}$} is closed \inserted{in $I_{\geq i}$,} because of $(*)$, so $J$ is indeed locally closed if we have $(*)$, the converse being even more obvious \ldots)

\uline{NB} One can also say, \inserted{if $\unsure{\mathscr{L}|\mathcal{L}}$ finite,} that the $A^\circ$ $(A\in\mathscr{L})$ are the \deleted{\unsure{at}} fibers of the Boolean algebra generated by the $A\in\mathscr{L}$ in $\unsure{\mathcal{P}|\mathscr{P}}(S)$ $(S=|\mathscr{L}|)$. The sets obtained from the $A\in\mathscr{L}$, by unions, intersections and \unsure{complements|complementations} in $S=\unsure{\lvert\mathscr{L})|\lvert\mathscr{L}\rvert}$, are exactly those saturated under the equivalence relation in $S$ defined by the $\unsure{A^\circ|X^\circ}$ \unsure{i.e.|id.} the one defined by $\varphi$. Thus they correspond bijectively to arbitrary subsets of $I$, by the form
\[
 J\longmapsto A_J=\varphi^{-1}(J)=\bigcup_{i\in\unsure{J|\supset}}A_i^\circ .
\]
We \unsure{shall call|shall call} \unsure{its|these|the} sets \unsure{thus|also} the $\mathscr{L}$-\uline{tame} or $\mathscr{L}$-\uline{constructible} \uline{subsets} of $\mathscr{L}$.
